\begin{explanationbox}
\textbf{\textcolor{CExplanation}{一句话直觉}}

等比数列中，每一项是公比的幂次，下标和对应指数和，乘积对应指数相加。

\medskip
\textbf{\textcolor{CExplanation}{核心拆解}}
\begin{description}[style=nextline, leftmargin=2.8em, labelsep=0.5em, itemsep=0.45em]
    \item[\textbf{要点一（比例关系）：}] 等比数列通项 $a_n = a_1 q^{n-1}$，两项相乘时指数相加：$a_m a_n = a_1^2 q^{m+n-2}$。
    \item[\textbf{要点二（指数和相等）：}] 若 $m+n=p+q$，则指数 $(m+n-2)=(p+q-2)$，因此乘积相等。
\end{description}

\medskip
\textbf{\textcolor{CExplanation}{几何本质（指数对称性）}}

将下标视为指数函数的自变量，乘积对应指数加法群中的运算。

\medskip
\textbf{\textcolor{CExplanation}{代数意义（幂运算性质）}}

$a_m a_n = a_1^2 q^{m+n-2}$ 只依赖于下标和，与具体下标无关。

\medskip
\textbf{\textcolor{CExplanation}{考点价值（简化计算）}}
\begin{itemize}[leftmargin=*, itemsep=0.5em, topsep=0.4em]
    \item \textbf{考法一（求乘积）：} 已知两项求另两项乘积，无需通项。
    \item \textbf{考法二（证明对称性）：} 用于证明等比数列中若干等式。
\end{itemize}

\medskip
\textbf{\textcolor{CExplanation}{顿悟点}}

乘积结果只与下标和有关，与具体项位置无关。

\medskip
\textbf{\textcolor{CExplanation}{使用场景}}
\begin{itemize}[leftmargin=*, itemsep=0.5em, topsep=0.4em]
    \item \textbf{场景一（求特定项）：} 已知 $a_2 a_8=16$，求 $a_5$。
    \item \textbf{场景二（化简条件）：} 已知 $a_1 a_3 = a_2^2$ 可直接推出公比关系。
\end{itemize}
\end{explanationbox}
